\documentclass[10pt]{article}
\usepackage[utf8]{inputenc}
\usepackage[T1]{fontenc}
\usepackage{amsmath}
\usepackage{amsfonts}
\usepackage{amssymb}
\usepackage[version=4]{mhchem}
\usepackage{stmaryrd}
\usepackage{graphicx}
\usepackage[export]{adjustbox}
\graphicspath{ {./images/} }
\usepackage{bbold}

\DeclareUnicodeCharacter{27F9}{\ifmmode\Longrightarrow\else{$\Longrightarrow$}\fi}

\begin{document}
\section*{Some basic game theory}
ECN 714\\
$2$

Bayesian games

\section*{Introduction}
\begin{itemize}
  \item In many games of economic interest, players may have imprecise information about some aspects of the game.
  \item For instance, a player may not be certain about her oponents' exact payoffs.
  \item A "representation" of the strategic situation described is necessary.
\end{itemize}

\section*{Example}
\begin{itemize}
  \item Players 1 and 2 choose actions simultaneously and independently; Player 1 chooses Up or Down, player 2 chooses Left or Right.
  \item Player 2 knows the game's payoffs.
\end{itemize}

\begin{center}
\begin{tabular}{lll}
 & $L$ & $\mathbf{R}$ \\
\hline
$U$ & 1,2 & 0,1 \\
\hline
D & 0,4 & 1,3 \\
\hline
\end{tabular}
\end{center}

\begin{itemize}
  \item Player 1 knows her own payoffs but she is not certain about 2's payoffs.
\end{itemize}

\section*{Player 1 as a Bayesian player}
\begin{itemize}
  \item Player 1 knows his own payoffs but she is not certain about 2's payoffs.
\end{itemize}

$$
\begin{array}{ccc} 
& \mathbf{L} & \mathbf{R} \\
\hline \mathrm{U} & 1 & 0 \\
\hline \mathrm{D} & 0 & 1
\end{array}
$$

\begin{itemize}
  \item Player 1 believes that player 2's payoffs are either
\end{itemize}

\begin{center}
\begin{tabular}{lllllll}
 & $\mathbf{L}$ & $\mathbf{R}$ &  &  & $\mathbf{L}$ & $\mathbf{R}$ \\
\hline
U & 2 & 1 & or & U & 3 & 4 \\
\hline
D & 4 & 3 &  & D & 1 & 2 \\
\hline
\end{tabular}
\end{center}

and assigns probability 0.6 to the first alternative and 0.4 to the second one.

\begin{itemize}
  \item Combine payoffs in a table, let states be $t=a$ and $t=b$ and player 1's beliefs are $p_{a}=0.6$ and $p_{b}=0.4$.
\end{itemize}

\begin{center}
\begin{tabular}{ccc}
\multicolumn{3}{c}{$t=a, p_{a}=.6$} \\
\hline
 & $L$ & R \\
\hline
$U$ & 1,2 & 0,1 \\
\hline
D & 0,4 & 1,3 \\
 & $t=b, p_{b}=.4$ &  \\
\hline
 & L & $R$ \\
\hline
U & 1,3 & 0,4 \\
\hline
$D$ & 0,1 & 1,2 \\
\hline
\end{tabular}
\end{center}

\section*{Benchmark: complete information}
\begin{itemize}
  \item Complete information: both players know whether $t=a$ or $t=b$.
  \item Strategies are functions:\\
$\left[s_{1}(a)=U, s_{1}(b)=D\right],\left[s_{2}(a)=L, s_{2}(b)=R\right]$
  \item Solution:
\end{itemize}

If $t=a$ then a NE is $(U, L)$.\\
If $t=b$ then a NE is $(D, R)$.

\begin{itemize}
  \item If listed probabilities are correct, an outside observer will see $60 \%$ of the time ( $U, L$ ) and 40\% of the time ( $D, R$ ).
\end{itemize}

\section*{Incomplete information}
\begin{itemize}
  \item Player 1 does not observe $t$; player 2 does. From 1's viewpoint, 2's strategies are functions.
  \item If $t=a, 2$ plays $L$ because $L$ is dominant.
  \item If $t=b, 2$ plays $R$ because $R$ is dominant.
  \item Player 1 has two strategies $U$ and $D$. Single information set.
\end{itemize}

$$
\begin{aligned}
u^{1}(U) & =0.6 \times 1+0.4 \times 0=0.6 \\
u^{1}(D) & =0.6 \times 0+0.4 \times 1=0.4
\end{aligned}
$$

Hence, player 1 chooses $U$.

\section*{Conclusions}
\begin{itemize}
  \item The intuitive outcome of the game was identified.
  \item Under incomplete information, since player 1 does not know the state $t$, she can no longer change her action depending on it.
  \item An outside observer would see player 1 choosing $U$ every time, and player 2 changing so that 60\% of time he plays $L$ and 40\% of the time he plays $R$.
  \item The incomplete information game is not a combination of the two complete information games.
\end{itemize}

\section*{Bayesian games}
\begin{itemize}
  \item Goal: to obtain a game representation and a solution concept to generalize the example.
  \item Incomplete information is modeled by adding a move by nature before the game begins to the extensive form representation.\\
Nature thus determines the realization of players' private information. Harsanyi (1967), Harsanyi (1968a), Harsanyi (1968b).
  \item Uncertainty about payoffs or other elements is a move by nature that takes place at the root of the game.
  \item To model a player that is not certain about its opponents, a Bayesian game creates types for the opponents (not for the player that lacks the information.) Each type represents a possible realization of the information.
  \item A Bayesian game is the normal form representation of the game played before the private information is revealed.
  \item A Bayesian Nash equilibrium (BNE) is a Nash equilibrium of the normal form game described.
\end{itemize}

\section*{Harsanyi's idea applied to the example}
\includegraphics[max width=\textwidth, alt={}, center]{07361299-5ece-411c-bdba-048b086fbf1e-13_1393_1264_348_175}\\
\includegraphics[max width=\textwidth, alt={}, center]{07361299-5ece-411c-bdba-048b086fbf1e-13_1396_1279_348_1463}

\section*{Bayesian games-definitions}
Definition 1 (Private information) Private information is the information that an agent has that is not available to any other agent.

\begin{itemize}
  \item Intuitively in a Bayesian game, players first observe their private information or type.
  \item Players have beliefs about opponents' types. In the example, player 1's beliefs are that player 2 has payoffs "a" with probability 0.6 and payoffs "b" with probability 0.4 .
  \item Each player simultaneously and idependently chooses his or her action.
  \item Final payoffs are determined by realization of information and actions chosen.
\end{itemize}

Definition 2 (Bayesian game) A Bayesian game $\Gamma$\\
$=\left[\mathcal{J},\left\{T_{i}, A_{i}, U_{i}, \beta_{i}\right\}_{i \in \mathcal{J}}\right]$ is a collection of

\begin{enumerate}
  \item $\mathcal{J}$, finite set of players.
  \item $T_{i}$, set of player $i$ 's private information or types.
  \item $A_{i}$, set of actions available to player $i$.
  \item $U_{i}: T \times A \longrightarrow \mathbb{R}$, player $i$ 's payoff function, where $T=\prod_{i \in \mathcal{J}} T_{i}$ and $A=\prod_{i \in \mathcal{J}} A_{i}$.
  \item $\beta_{i}: T_{i} \rightarrow M\left(T_{-i}\right)$, player $i$ 's beliefs about the other players' private information.\\
$M\left(T_{-i}\right)$ is set of probability distritubions over $T_{-i}$.
\end{enumerate}

\begin{itemize}
  \item When $T_{-i}$ is finite, $\beta_{i}\left(t_{-i} \mid t_{i}\right)$ is the probability that $i$ assigns to $t_{-i}$ when $i$ 's information is $t_{i}$.
  \item More generally, $\beta_{i}\left(\cdot \mid t_{i}\right)$ is a probability distribution on $T_{-i}$ for each $t_{i}$.
\end{itemize}

\section*{Beliefs}
Definition 3 (Consistent beliefs) Beliefs are consistent with a probability $P$ on $T$ if for all $i \in N$, and for all $\left(t_{-i}, t_{i}\right) \in T$,

$$
\beta_{i}\left(t_{-i} \mid t_{i}\right)=\frac{P\left(t_{-i}, t_{i}\right)}{\sum_{t_{-i} \in T_{-i}} P\left(t_{-i}, t_{i}\right)} .
$$

Definition 4 (Common prior) When beliefs are consistent with a given probability distribution $P$, it is said that players have a common prior $P$.

\section*{Beliefs and common prior}
\begin{center}
\includegraphics[max width=\textwidth, alt={}]{07361299-5ece-411c-bdba-048b086fbf1e-18_1550_2286_354_163}
\end{center}

Two players, $T_{1}=\{i, u\}$, say informed or uninformed; and $T_{2}=\{a, b\}$. Then,

$$
\begin{array}{llll}
P\left(t_{1}, t_{2}\right) & t_{2}=a & t_{2}=b & \sum_{t_{2} \in T_{2}} \\
\hline t_{1}=i & P(i, a) & P(i, b) & P_{T_{1}}(i) \\
\hline t_{1}=u & P(u, a) & P(u, b) & P_{T_{1}}(u) \\
\hline \sum_{t_{1} \in T_{1}} & P_{T_{2}}(a) & P_{T_{2}}(b) & 1 \\
& \sum_{t_{1} \in T_{1}, t_{2} \in T_{2}} P\left(t_{1}, t_{2}\right)=1 &
\end{array}
$$

\section*{Example: Consistent beliefs}
\begin{center}
\begin{tabular}{llll}
 & a & b & $\sum_{a, b}$ \\
\hline
i & 0 & 0 & 0 \\
\hline
u & 0.6 & 0.4 & 1 \\
\hline
$\sum_{i, u}$ & 0.6 & 0.4 &  \\
\hline
\end{tabular}
\end{center}

$$
\begin{aligned}
& \beta_{2}\left(t_{1}=i \mid a\right)=\frac{P(i, a)}{P(i, a)+P(u, a)}=\frac{0}{0+.6}=0 \\
& \beta_{2}\left(t_{1}=i \mid b\right)=\frac{P(i, b)}{P(i, b)+P(u, b)}=\frac{0}{0+.4}=0
\end{aligned}
$$

$\beta_{1}\left(t_{2}=a \mid i\right)=$ not defined by Bayes rule.

$$
\begin{aligned}
& \beta_{1}\left(t_{2}=a \mid u\right)=\frac{P(u, a)}{P(u, a)+P(u, b)}=\frac{.6}{1}=.6 \\
& \beta_{1}\left(t_{2}=b \mid u\right)=\frac{P(u, b)}{P(u, a)+P(u, b)}=\frac{.4}{1}=.4
\end{aligned}
$$

\section*{Beliefs that are not consistent}
\begin{center}
\begin{tabular}{llll}
 & $\mathbf{a}$ & $\mathbf{b}$ & $\sum_{a, b}$ \\
\hline
i & 0.06 & 0.04 & 0.10 \\
\hline
u & 0.54 & 0.36 & .90 \\
\hline
$\sum_{i, u}$ & 0.6 & 0.40 &  \\
\hline
\end{tabular}
\end{center}

$$
\begin{aligned}
\beta_{1}\left(t_{2}=a \mid i\right) & =0.6 \\
\beta_{1}\left(t_{2}=b \mid i\right) & =0.4 \\
\beta_{1}\left(t_{2}=a \mid u\right) & =0.54 / 0.9=0.6 \\
\beta_{1}\left(t_{2}=b \mid u\right) & =0.36 / 0.9=0.4
\end{aligned}
$$

\section*{Strategies}
Fix a Bayesian game $\Gamma$.

\begin{itemize}
  \item A pure strategy $s_{i}$ for player $i$ is any function $s_{i}: T_{i} \rightarrow A_{i}$. The set of player $i$ 's pure strategies is denoted by $S_{i}$.
  \item A mixed strategy for player $i$ is a probability distribution $\mu_{i} \in M\left(S_{i}\right)$, the set of probability distributions on $S_{i}$.
  \item A behavior strategy for player $i$ is any function $\sigma_{i}: T_{i} \rightarrow M\left(A_{i}\right)$. The set of player $i$ 's behavior strategies is denoted by $\Sigma_{i}$.
  \item A strategy profile a strategy profile is a collection of strategies, one for each player in the game.
\end{itemize}

\section*{Observations}
Let $\Gamma$ be a Bayesian game with two players.\\
A pure-strategy profile $\left(s_{1}^{\prime}, s_{2}^{\prime}\right)$ is a Bayesian Nash equilibrum (BNE) if ( $\forall i \in \mathcal{J})\left(\forall t_{i} \in T_{i}\right)\left(\forall s_{i} \in S_{i}\right)(\forall j \neq i)$,


\begin{align*}
& \sum_{t_{j} \in T_{j}} u_{i}\left(t_{j}, t_{i}, s_{j}^{\prime}\left(t_{j}\right), s_{i}^{\prime}\left(t_{i}\right)\right) \beta_{i}\left(t_{j} \mid t_{i}\right) \geq \\
& \quad \sum_{t_{j} \in T_{j}} u_{i}\left(t_{j}, t_{i}, s_{j}^{\prime}\left(t_{j}\right), s_{i}\left(t_{i}\right)\right) \beta_{i}\left(t_{j} \mid t_{i}\right) . \tag{1}
\end{align*}


\begin{itemize}
  \item Note that $s_{i}\left(t_{i}\right)$ could be any element of $A_{i}$
\end{itemize}

\section*{Observations 2}
Thus equivalently,

$$
\begin{aligned}
\sum_{t_{j} \in T_{j}} u_{i}\left(t_{j}, t_{i},\right. & \left.s_{j}^{\prime}\left(t_{j}\right), s_{i}^{\prime}\left(t_{i}\right)\right) \beta_{i}\left(t_{j} \mid t_{i}\right) \geq \\
& \sum_{t_{j} \in T_{j}} u_{i}\left(t_{j}, t_{i}, s_{j}^{\prime}\left(t_{j}\right), a_{i}\right) \beta_{i}\left(t_{j} \mid t_{i}\right), \forall a_{i} \in A_{i}
\end{aligned}
$$

\begin{itemize}
  \item Notation: Using expectations
\end{itemize}

$$
E\left[u_{i} \mid t_{i}, s_{i}^{\prime}, s_{j}^{\prime}, \beta_{i}\right]=\sum_{t_{j} \in T_{j}} u_{i}\left(t_{j}, t_{i}, s_{j}^{\prime}\left(t_{j}\right), s_{i}^{\prime}\left(t_{i}\right)\right) \beta_{i}\left(t_{j} \mid t_{i}\right) .
$$

\section*{Observations 3}
Let

$$
\begin{aligned}
E\left[u_{i} \mid s_{i}^{\prime}, s_{j}^{\prime}\right]=\sum_{t_{i} \in T_{i}} \sum_{t_{j} \in T_{j}} u_{i}\left(t_{j}, t_{i}, s_{j}^{\prime}\left(t_{j}\right), s_{i}^{\prime}\left(t_{i}\right)\right) & \\
\beta_{i}\left(t_{j} \mid t_{i}\right) P_{T_{i}}\left(t_{i}\right) &
\end{aligned}
$$

where $P_{T_{i}}$ is the marginal of $P$ on $T_{i}$.

\begin{itemize}
  \item Then, $\left(s_{i}^{\prime}, s_{j}^{\prime}\right)$ is a BNE if
\end{itemize}


\begin{equation*}
E\left[u_{i} \mid s_{i}^{\prime}, s_{j}^{\prime}\right] \geq E\left[u_{i} \mid s_{i}, s_{j}^{\prime}\right], \forall s_{i} \in S_{i} . \tag{2}
\end{equation*}


\begin{itemize}
  \item The inequality in Equation 1 implies the inequality in Equation 2. When $P_{T_{i}}\left(t_{i}\right)>0$ for every $t_{i}$, both inequalities are equivalent.
\end{itemize}

\section*{Dominant strategy}
\begin{itemize}
  \item A pure strategy $s_{i}$ is (weakly) dominant for player $i$ in $\Gamma$, if
\end{itemize}

$$
\forall s_{-i}, E\left[U_{i} \mid s_{-i}, s_{i}\right] \geq E\left[U_{i} \mid s_{-i}, s_{i}^{\prime}\right], \forall s_{i}^{\prime}
$$

\begin{itemize}
  \item A profile $s=\left\{s_{i}\right\}_{i=1}^{I}$ is a (weakly) dominant strategy equilibrium if $s_{i}$ is (weakly) dominant for $i=1, \ldots, n$.
\end{itemize}

\section*{Baysian Nash equilibrium}
Definition 5 (Bayesian Nash equilibrium) A behavior strategy profile $\sigma^{\prime} =\left(\sigma_{1}^{\prime}, \ldots, \sigma_{N}^{\prime}\right)$ is a Bayesian Nash equilibrum (BNE) if ( $\left.\forall i \in \mathcal{J}\right)\left(\forall t_{i}\right. \left.\in T_{i}\right)\left(\forall \sigma_{i} \in \Sigma_{i}\right)$,

$$
E\left[u_{i} \mid t_{i}, \sigma_{i}^{\prime}, \sigma_{-i}^{\prime}, \beta_{i}\right] \geq E\left[u_{i} \mid t_{i}, \sigma_{i}, \sigma_{-i}^{\prime}, \beta_{i}\right] .
$$

Note that

$$
\begin{aligned}
E\left[U_{i} \mid t_{i}, \sigma_{i}, \sigma_{-i}^{\prime}, \beta_{i}\right]= & \sum_{t_{-i} \in T_{-i}} \sum_{a_{1} \in A_{1}} \ldots \sum_{a_{I} \in A_{I}} u_{i}\left(t_{-i}, t_{i}, a_{1}, \ldots, a_{I}\right) \\
& \left(\prod_{j \neq i} \sigma_{j}^{\prime}\left(a_{j} \mid t_{j}\right)\right) \sigma_{i}\left(a_{i} \mid t_{i}\right) \beta_{i}\left(t_{-i} \mid t_{i}\right) .
\end{aligned}
$$

\section*{Universal type space}
Difficulty with incomplete information. Refer to leading example.

\begin{itemize}
  \item Player 2 observes $t_{2}$, player 1 does not. Player 1 must have beliefs about $t_{2}$.
  \item Depending on 1's beliefs, 2's preferred strategy may change.
  \item In the example, 2 has a dominant strategy in each game where a game is defined for each realization of $t_{2}$. So this is not an issue. It is in general.
  \item Player 2 does not observe 1's beliefs; 2 must form beliefs about player 1's beliefs.
  \item Player 1 optimal action may depend on player 2's choices. In turn these choices may depend on player 2's beliefs. Thus player 1 must form beliefs, about player 2's beliefs about player 1's beliefs.
  \item Continue ad infinitum.
\end{itemize}

See Brandenburger and Dekel (1993)\} and Mertens and Zamir (1985).

\section*{Example: auction, stylized representation}
\begin{center}
\includegraphics[max width=\textwidth, alt={}]{07361299-5ece-411c-bdba-048b086fbf1e-32_1559_2289_351_163}
\end{center}

\section*{Independent private values model (IPVM)}
A standard model in trading environments.

\begin{itemize}
  \item A seller has a single, indivisible object.
  \item There are $I$ potential buyers or bidders.
  \item Bidder $i$ privately observes her valuation $x_{i} \in[0,1]$ for the object.
  \item Preferences: If $i$ receives the object and pays an amount $m, i$ 's payoff is $x_{i}-m$. Expected utility is assumed.
  \item Beliefs: Variables $x_{i}$ are i.i.d.. Each $x_{i}$ has density $f$, assumed to be strictly positive and with full support. This is also a symmetry assumption, i.e. bidders are ex ante equal.
\end{itemize}

\section*{Notation}
\begin{itemize}
  \item $x=\left(x_{1}, \ldots, x_{I}\right), x_{-i}=\left(x_{1}, \ldots, x_{i-1}, x_{i+1}, x_{I}\right)$.
  \item $X_{1}, X_{2}, \ldots, X_{I}$ sets, $X=\prod_{i=1}^{I} X$.
  \item $M\left(X_{i}\right)$ is the set of probability distributions on $X_{i}$.
  \item First order statistic: $x^{(1)}=\max \left\{x_{i}\right\}_{i=1}^{I}$;
  \item Second order statistic: $x^{(2)}$
\end{itemize}

\section*{Classic auctions}
\begin{itemize}
  \item Sealed-bid, second-price auction (SPA).
  \item Bidders simultaneously and independently submit their bids in sealed envelopes.
  \item Envelopes are collected and opened.
  \item Highest bidder, the winner, is assigned the object. Ties in determining the highest bidder are resolved according to some predetermined procedure.
  \item The winner pays the second-highest bid.
  \item Strategies?
\end{itemize}

\section*{Second-price, sealed-bid auction (SPA)}
Theorem 1 (SPA) IPVM. The strategy profile $\left\{b_{i}\left(x_{i}\right)=x_{i}\right\}_{i=1}^{I}$ is a weakly dominant-strategy equilibrium (WDSE) of the SPA.

Proof. Let $b_{-i}$ be the vector of $i$ 's oponents bids.

$$
u_{i}\left(x_{i}, b_{i}, b_{-i}^{(1)}\right)= \begin{cases}0 & \text { if } b_{i}<b_{-i}^{(1)} \\ \left(x_{i}-b_{-i}^{(1)}\right) & \text { if } b_{i}>b_{-i}^{(1)} \\ \left(x_{i}-b_{-i}^{(1)}\right) P(\text { tie }) & \text { if } b_{i}=b_{-i}^{(1)}\end{cases}
$$

\begin{itemize}
  \item $x_{i}<b_{-i}^{(1)} \Longrightarrow i^{\prime} s$ best response is to lose, thus $b_{i}=x_{i}<b_{-i}^{(1)}$.
  \item $x_{i}>b_{-i}^{(1)} \Longrightarrow i^{\prime} s$ best response is to win, thus $b_{i}=x_{i}>b_{-i}^{(1)}$.
  \item $x_{i}=b_{-i}^{(1)} \Longrightarrow i$ is indifferent, thus $b_{i}=x_{i}$. QED
\end{itemize}

\section*{SPA-Remarks}
\begin{itemize}
  \item Weakly dominant strategy equilibrium (WDSE) means a strong recommendation.
\end{itemize}

Principle: winner's payment does not depend on winner's bid.

\begin{itemize}
  \item Result is stronger than WDSE of the implicit game. (Continuum technicality.)
  \item Where is independence used? And symmetry?
  \item Are there other BNE?
  \item Solution is efficient.
\end{itemize}

\section*{Classic auctions}
\begin{itemize}
  \item Sealed-bid second-price auction (SPA)
  \item Sealed-bid first-price auction (FPA): Bidders submit their bids in sealed envelopes simultaneously and independently. Envelopes are opened. Highest bidder gets the object and pays his or her own bid.
  \item English auction
  \item Dutch auction
\end{itemize}

\section*{Existence of BNE}
Finite Bayesian games. Existence of BNE is a consequence of the general existence of NE in finite games.

Existence proofs use a fixed point argument.

\begin{itemize}
  \item Primitive payoff function: $U_{i}\left(x_{1}, x_{2}, a_{1}, a_{2}\right)$
  \item Behavior strategy: $\sigma_{i}: X_{i} \rightarrow M\left(A_{i}\right)$
  \item Best response correspondence:\\
$\psi_{i}\left(\sigma_{-i}\right)=\arg \max _{\sigma_{i}} E\left[U_{i} \mid \sigma_{i}, \sigma_{-i}\right]$
  \item $\psi(\sigma)=\prod_{i=1}^{I} \psi_{i}\left(\sigma_{-i}\right)$.
  \item Theorem of the Maximum (Berge): Continuity of primitive function implies upper-hemicontinuity of the argmax.\\
In this case, continuity of $E\left[U_{i} \mid \sigma_{1}, \sigma_{2}\right]$ implies uhc of $\psi$.\\
See for instance, Aliprantis and Border (2006), page 569.
  \item uhc and convex-valued $\psi$ plus compactness of domain ⟹ fixed point, therefore existence.
\end{itemize}

\section*{BNE in infinite games}
\begin{itemize}
  \item No general existence result for BNE in infinite games, games where type or action spaces have the cardinality of the continuum.
  \item Advantages of continuum model
  \item Modeling. Perhaps more appropriate, bids and uncertainty. Bids in practice are discrete. Should results depend on whether permitted bids are in pennies or nickels?
  \item Technical.
  \item Avoid complications of multiplicity of equilibria from mixing between discrete alternatives.
  \item Calculus.
\end{itemize}

\section*{BNE in infinite games}
\begin{itemize}
  \item Difficulty with continuum model.
  \item $\left.E\left[U_{i} \mid \sigma_{1}, \sigma_{2}\right)\right]=\int_{X_{1} \times X_{2}} \int_{A_{1}} \int_{A_{2}} U_{i}(x, a) \mathrm{d} \sigma_{1}\left(x_{1}\right) \mathrm{d} \sigma_{2}\left(x_{2}\right) \mathrm{d} \mu$
  \item Topology (metric) that makes $\psi$ uhc, has too many open sets to make the domain of $\psi$ compact and viceversa.
  \item Additional questions.
  \item Existence of pure strategy equilibria or purification of mixed strategy equilibria.
  \item Approximation discrete-continuum.
\end{itemize}

\section*{References}
Aliprantis, Charalambos D., and Kim C. Border. 2006. Infinite Dimensional Analysis, a Hitchhiker's Guide. Third edition. Berlin-Heidelberg: Springer-Verlag.\\
Brandenburger, Adam, and Eddie Dekel. 1993. "Hierarchies of Beliefs and Common Knowledge." Journal of Economic Theory 59: 189-98.\\
Harsanyi, John. 1967. "Games with Incomplete Information Played by Bayesian Players. Part i: The Basic Model." Management Science 14: 159-82.\\
---. 1968a. "Games with Incomplete Information Played by Bayesian Players. Part II: Bayesian Equilibrium Points." Management Science 14: 320-34.\\
---. 1968b. "Games with Incomplete Information Played by Bayesian Players. Part III: The Basic Probability Distribution of the Game." Management Science 14: 486-502.\\
Mertens, J.-F., and S. Zamir. 1985. "Formulation of Bayesian Analysis for Games with Incomplete Information." International Journal of Game Theory 10: 619-32.


\end{document}